This chapter presents the primary research contributions of this dissertation research: three sonic sculptures that function as research instruments. These works are titled \emph{Buffalo Soldier}, \emph{Did Sun Ra Fear Cryosleep?}, and \emph{Don't Play with My Hair}. I find it is important to articulate that these works should not be seen as illustrations of theoretical concepts developed elsewhere, nor should they be seen as artistic projects with research retroactively applied. They \textit{are} the research. Each work enacts complete cycles (or at times partial cycles or individual procedures) of Speculocultural Technopoiesis (ST), generating knowledge through the embodied process of making, modifying, and materializing culturally grounded technological practice.

The title of this chapter, "Sonic Speculocultural Technopoiesis," intentionally mirrors the dissertation's title itself, emphasizing that the practice and the research are inseparable. In research-creation, as Loveless argues, inquiry follows "driven curiosity as its lure and guide," allowing practices to emerge responsively rather than conform to predetermined frameworks \parencite[105]{Loveless_2019}. As a Black scholar-practitioner working with digital signal technologies, physical computing platforms, and digital signal processing software, these works emerged from my specific encounters with tools that often carry embedded assumptions about who makes, how we make, and what counts as legitimate technological practice. They represent sites where cultural knowledge and technical constraint friction into new procedural possibilities.

I call these works \textit{sonic sculptures} because they exist at the intersection of sound, materiality, spatiality, and speculation. They are physical objects embedded with sensors and computational systems. They are software patches and modified algorithms. They are speculative narratives that refuse dominant historical and technological paradigms. They are interactive experiences that demand embodied participation. Most importantly, they are experimental, their final forms remain emergent, responsive to what the practice itself reveals through iterative cycles of unsettling, remixing, embodying, and transmitting.

Each work is presented here through the lens of ST's cycle and domains. I describe the \textit{desire} or conceptual spark that initiated each work (often arising from personal lived experience and cultural investigation), the \textit{signal} or material and sonic form each work takes (the substrates, sensors, sounds, and software modifications), and the \textit{intention} or what each work seeks to accomplish culturally and technologically. Because these are experimental works developed through embodied practice, I must also acknowledge what remains uncertain: directions they may take, modifications still under development, and frictions not yet resolved. While each work stands as a distinct investigation, they inform one another iteratively. The technical insights from one shapes the approaches in another, conceptual frictions in one opening possibilities in the next.

\begin{comment}
This chapter demonstrates that I have a practice, that this practice is rigorous and accountable to specific cultural epistemologies, and that this practice generates transmissible knowledge through the creation of these three research instruments. The works are presented in the chronological order of their initial conceptual spark: \emph{Buffalo Soldier}, \emph{Did Sun Ra Fear Cryosleep?}, and \emph{Don't Play with My Hair}. While each work stands as a distinct investigation, they inform one another iteratively. The technical insights from one shapes the approaches in another, conceptual frictions in one opening possibilities in the next. Together, they constitute a body of practice-based research that attempts to refuse the separation between making and knowing, between cultural speculation and technological modification, between artistic creation and scholarly contribution.
\end{comment}